% !PRINTED 165
\weberpage{165}
%<<CONT
and $C'$. From this we conclude that $C$ also belongs to the invariants
of the group $S$, and state this as a theorem:

4. \gesperrt{If from an invariant form of the
group $S$ one forms arbitrary covariants, one obtains
new invariant forms of the group.}

\begin{center}
\S.\ 41.

\gesperrt{Hilbert's Theorem.}
\end{center}

The proof of the theorem on the finiteness of the system of invariants
 of a linear substitution group rests on a very
general theorem about systems of forms of any kind whatever, which
\gesperrt{Hilbert} discovered and applied with distinguished
success in manifold investigations on the finiteness
of systems of invariants, and to whose derivation we now wish
to proceed $^{1}$).

I. \gesperrt{If $\mathfrak{S}$ denotes any system of forms in the
$n$ variables $x_1,x_2,\ldots x_n$, in finite or
infinite number, then from $\mathfrak{S}$ a
finite number of forms $F_1,F_2,\ldots,F_\mu$ can be so
selected that every form $F$ of $\mathfrak{S}$ can be represented
by an expression}
\weq{1}{F = A_1F_1 + A_2F_2 + \cdots + A_\mu F_\mu}
\gesperrt{in which $A_1,A_2,\ldots A_\mu$ are
forms of the variables $x_1,x_2,\ldots x_n$.}

The definition of the system of forms $\mathfrak{S}$ must be so complete
that of every single form of the variables $x$ it is decided
whether it belongs to $\mathfrak{S}$ or not, but it is otherwise bound to no
assumption.

If $\mathfrak{S}$ consists of only a finite number of forms, then
our theorem is self-evident, for in this case one can
take all the forms of $\mathfrak{S}$ for $F_1,F_2,\ldots F_\mu$. The
proof will therefore only have to concern itself with the case of an infinite
system $\mathfrak{S}$.

For simplicity of expression we shall call the system of forms
 $F_1,F_2,\ldots F_\mu$ a \gesperrt{basis of the system} $\mathfrak{S}$.
The functions $A_1,A_2,\ldots A_\mu$ must be so constituted that

\bigskip
\noindent\rule{4cm}{0.4pt}

\noindent $^{1}$) \gesperrt{Hilbert}, ``Ueber die Theorie der algebraischen Formen''. Mathematische
 Annalen, Bd. 36 (1890).
%CONT>>
